\documentclass[11pt]{article}
\usepackage[margin=2.0cm]{geometry}
\usepackage{booktabs}
\usepackage{amsmath}
\usepackage{amssymb}
\usepackage{enumitem}
\usepackage{array}
\usepackage{listings}
\usepackage{xcolor}
\usepackage[T1]{fontenc}
\usepackage{lmodern}

\setlength{\parindent}{0pt}
\setlength{\parskip}{3.6pt}
\newcommand{\code}[1]{\texttt{#1}}
\newcommand{\hr}{\vspace{2pt}\hrule\vspace{6pt}}
\renewcommand{\arraystretch}{1.12}

\definecolor{kw}{RGB}{0,90,160}
\definecolor{cmt}{RGB}{110,110,110}
\lstdefinestyle{vhdl}{
  language=VHDL, basicstyle=\ttfamily\footnotesize,
  keywordstyle=\color{kw}\bfseries, commentstyle=\color{cmt}\itshape,
  showstringspaces=false, columns=fullflexible, breaklines=true,
  frame=leftline, framesep=6pt, xleftmargin=8pt, aboveskip=6pt, belowskip=6pt
}
\lstset{style=vhdl}

\begin{document}

\begin{center}
{\Large\bfseries EE\,362 --- Computer Organization Laboratory}\\[2pt]
{\large\bfseries Lab 8 --- Guided Examples: Solution}\\[2pt]
{\itshape Ex 1 absolute targets \quad$\cdot$\quad Ex 2 the condition equations \quad$\cdot$\quad Ex 3 two decisions, not one}\\[4pt]
Fall 2026 \quad$\cdot$\quad Lab TA: Gandhi Torres
\end{center}
\hr

\begin{center}
\small
\begin{tabular}{cl>{\raggedright\arraybackslash}p{6.0cm}>{\raggedright\arraybackslash}p{3.8cm}}
\toprule
& circuit & new idea & feeds the assignment as \\
\midrule
1 & \code{absolute\_target} & a jump target is \textbf{wiring}, not arithmetic & \code{NIAG\_Abslt} \\
2 & \code{branch\_conditions} & six conditions from \textbf{two flags} & \code{NIAG\_Cond} \\
3 & \code{pc\_select} & \code{Sel} and \code{Cond} are \textbf{separate} & \code{NIAG\_MUX} \\
\bottomrule
\end{tabular}
\end{center}

\vspace{4pt}\hr

\section*{Example 1 --- The absolute target is pure wiring}

\begin{lstlisting}
Abslt <= NSA(31 downto 28) & Index & "00";
\end{lstlisting}

\begin{center}
\begin{tabular}{lll}
\toprule
bits & source & why \\
\midrule
4 & \code{NSA(31..28)} & \textbf{inherited} --- the instruction cannot supply them \\
26 & \code{Index} & all the room the encoding has \\
2 & \code{"00"} & instructions are word-aligned \\
\bottomrule
\end{tabular}
\end{center}

\textbf{There is no adder in this component.} Concatenation is free: it is a
renaming of wires, and it costs zero logic in the Fitter report.

\subsection*{Worked example from the Lab 12 program}

\code{J 0x08} at \code{0x14}, encoded \code{0x08000002}:
\code{NSA} $=$ \code{0x18}, top four bits \code{0000}, \code{Index} $= 2$, so
\code{Abslt} $=$ \code{0x00000008}.

\subsection*{The consequence: a jump cannot leave its region}

Because the top four bits are inherited, a \code{J} reaches only addresses in the
\textbf{same 256\,MB region} ($2^{28}$ bytes) as \code{NSA}. To go further you
need \code{JR}, whose address comes from a register and therefore supplies all 32
bits.

\vspace{4pt}\hr

\section*{Example 2 --- Six conditions from two flags}

\begin{center}
\begin{tabular}{lll}
\toprule
code & meaning & Taken \\
\midrule
\code{NONE} & always & $1$ \\
\code{EQ} & $=0$ & $Z$ \\
\code{NEQ} & $\neq 0$ & $\overline{Z}$ \\
\code{GTZ} & $>0$ & $\overline{N}\cdot\overline{Z}$ \\
\code{GEZ} & $\geq 0$ & $\overline{N}$ \\
\code{LTZ} & $<0$ & $N$ \\
\code{LEZ} & $\leq 0$ & $N + Z$ \\
\bottomrule
\end{tabular}
\end{center}

\begin{quote}
\textbf{Every conditional branch in the machine is one or two gates over the ALU's
two flags.} That is the whole conditional-control-flow mechanism.
\end{quote}

\subsection*{2.1 The complement pairs}

\[
\code{GTZ} \leftrightarrow \code{LEZ} \qquad
\code{GEZ} \leftrightarrow \code{LTZ} \qquad
\code{EQ}  \leftrightarrow \code{NEQ}
\]

By De Morgan,
$\overline{\code{GTZ}} = \overline{\overline{N}\cdot\overline{Z}} = N + Z = \code{LEZ}$.
So seven conditions do not cost seven circuits --- three gates and their
complements cover everything, and \code{Cond} merely selects which reaches the
output.

\subsection*{2.2 The two that need both flags}

\code{GTZ} and \code{LEZ} are the only conditions requiring \emph{both} $N$ and
$Z$, and they are where the mistakes happen:

\begin{center}
\begin{tabular}{lll}
\toprule
wrong version & agrees with correct except & symptom \\
\midrule
\code{GTZ} $= \overline{N}$ & at $N{=}0, Z{=}1$ (result exactly zero) & branches when the value is zero \\
\code{LEZ} $= N$ & at $N{=}0, Z{=}1$ & \textbf{fails to branch at zero} \\
\bottomrule
\end{tabular}
\end{center}

\textbf{The second one is the Lab 12 killer:} the reference program's loop ends
with \code{BLEZ}, so a \code{LEZ} that forgets $+Z$ runs the loop past its exit
and the CPU never terminates. Test $N{=}0, Z{=}1$ explicitly --- it is one vector.

\subsection*{2.3 Why the default is \emph{not taken}}

\begin{lstlisting}
'0' when others;        -- an unknown condition code does NOT branch
\end{lstlisting}

\begin{quote}
\textbf{A CPU that takes a branch it does not understand runs away; one that
ignores it merely does the wrong thing slowly.} Between two failures, choose the
one you can debug.
\end{quote}

That is a safety argument, and it is the same reasoning as \code{CTRL\_DEFAULT} in
Lab 10: unknown input, harmless behaviour.

\vspace{4pt}\hr

\section*{Example 3 --- Two decisions, not one}

Students routinely try to merge these into a single control signal. They are
genuinely two separate questions:

\begin{center}
\begin{tabular}{ll}
\toprule
\code{NIAG\_sel} & \textbf{which} address is the candidate? (\code{NSA} / \code{REL} / \code{REG} / \code{ABS}) \\
\code{NIAG\_cond} & is a \textbf{relative} branch actually taken? (from the ALU flags) \\
\bottomrule
\end{tabular}
\end{center}

\begin{center}
\begin{tabular}{ll}
\toprule
\code{Sel} & $PC^{+}$ \\
\midrule
\code{NSA} & \code{NSA} \hfill unconditional \\
\code{REL} & \textbf{\code{Taken} ? \code{Rel} : \code{NSA}} \hfill \textbf{conditional} \\
\code{REG} & \code{RegAddr} \hfill unconditional \\
\code{ABS} & \code{Abslt} \hfill unconditional \\
\bottomrule
\end{tabular}
\end{center}

\subsection*{3.1 The structure: a gate in front of a mux}

\begin{lstlisting}
rel_or_nsa <= Rel when Taken = '1' else NSA;     -- 2-to-1: the Taken gate

with Sel select
    NextPC <= NSA        when NIAG_SEL_NSA,
              rel_or_nsa when NIAG_SEL_REL,       -- only this path is conditional
              RegAddr    when NIAG_SEL_REG,
              Abslt      when NIAG_SEL_ABS,
              NSA        when others;             -- safe default
\end{lstlisting}

\textbf{This is your Lab 3 work again:} a 2-to-1 multiplexer feeding a 4-to-1
multiplexer. Nothing new is being built --- which is worth saying, because the
component is called ``next instruction address generation'' and sounds like it
should be hard.

\subsection*{3.2 Why not-taken falls back to \code{NSA}}

A branch that is not taken is not a special state. It means ``carry on to the next
instruction'' --- which is exactly \code{NSA}.

\begin{quote}
\textbf{That is why there is no fifth select code for ``branch not taken''.} The
fall-back is already there.
\end{quote}

\subsection*{3.3 Why jumps are unconditional by construction}

There is no conditional \code{J} in this instruction set. A jump's condition would
have to come from the same flags --- and that is what branches are for.

\begin{quote}
\textbf{A design decision, not a limitation:} conditional control flow uses
relative addressing because conditions are usually local, and unconditional
transfers use absolute addressing because they usually are not.
\end{quote}

\subsection*{3.4 The vector pair that proves you got it right}

\begin{center}
\begin{tabular}{lll}
\toprule
vectors & expected & proves \\
\midrule
\code{Sel}=\code{REL}, \code{Taken} $0 \to 1$ & \code{NSA}, then \code{Rel} & the branch path \emph{is} conditional \\
\code{Sel}=\code{ABS}, \code{Taken} $0 \to 1$ & \code{Abslt} both times & the jump path is \emph{not} \\
\bottomrule
\end{tabular}
\end{center}

If both pairs behave the same way, \code{Sel} and \code{Cond} have been merged
somewhere.

\vspace{4pt}\hr

\section*{Into the assignment}

\begin{center}
\begin{tabular}{lll}
\toprule
component & what it does & storage? \\
\midrule
\code{NIAG\_Abslt} & concatenation & no \\
\code{NIAG\_Cond} & six equations over $Z$, $N$ & no \\
\code{NIAG\_MUX} & \code{Taken} gate $+$ 4-to-1 select & no \\
\code{NIAG} & wires the whole unit together, including Lab 7's parts & \textbf{one register: the PC} \\
\bottomrule
\end{tabular}
\end{center}

\textbf{The whole address-generation unit contains exactly one storage element.}
Everything else is combinational, because the next address must be ready
\emph{within} the cycle that computes it. Check it in the Fitter report: with
\code{NIAG} as top level the register count is \textbf{32} --- the program counter
and nothing else.

\vspace{6pt}
\hr
{\small \textbf{Forward.} In Lab 10 the decoder sets \code{NIAG\_sel} and
\code{NIAG\_cond} for each instruction: \code{J} sets \code{ABS}, \code{BLEZ} sets
\code{REL} with \code{LEZ}, and everything else sets \code{NSA}. This unit does
not change again.}

\end{document}
